% sections/05_accounting.tex -- Section 5, Welfare and Fiscal Accounting (W2).
% Plan: BLUEPRINT sections 2.5 and 4 (equations eq:ce ... eq:bcmap, in that order), Overrides O1 and O5.
% Sources: dossier 03 section 3; frot/sub2025/12_app_quant.tex:563-640; 7_policy_reforms.tex:60-80, 117-121, 400-425.
% Every sentence with a number from a table, figure or source line carries a "% src:" comment (Override O5).
\section{Welfare and Fiscal Accounting}\label{sec:accounting}

Every reform in Sections~\ref{sec:elimination}--\ref{sec:menu} is evaluated with one welfare measure and one fiscal accounting. This section defines both, maps each column of the reform tables to a formula, and relates the marginal value of public funds to the objects of the sufficient-statistics literature; \autoref{app:derivations} collects the algebra. Superscript $0$ denotes the status quo, $\rho$ a reform, and $\Delta$ the difference.

\subsection{Welfare measures and willingness to pay}\label{sec:welfare}

Let $V^{i}_{30}$ be the value function \eqref{eq:bellman} of simulated individual $i$ at age 30, the entry age of the model.
% src: 5_quantitative_model.tex:23 (entry age 30)
Cohort welfare is the average of $V^{i}_{30}$ over the $I^{s}$ simulated individuals, $V^{0}$ under the status quo and $V^{\rho}$ under reform $\rho$: an ex ante measure, taken before health, mortality, wage, and medical-expense shocks are realized. The consumption-equivalent (CE) welfare change $x^{\CE}_{\rho}$ is the log change in status-quo consumption at every age and state, with leisure held fixed, that makes the cohort as well off as under the reform; $e^{x^{\CE}_{\rho}}-1$ is the corresponding proportional change, which $x^{\CE}_{\rho}$ approximates to first order. Flow utility \eqref{eq:utility} is homogeneous of degree $\gamma(1-\nu)$ in consumption, so $x^{\CE}_{\rho}$ solves $V^{\rho}=e^{\gamma x^{\CE}_{\rho}(1-\nu)}V^{0}$:
% src: 12_app_quant.tex:567-570 (CE change defined through V_after = e^{gamma x (1-nu)} V_before)
\begin{equation}\label{eq:ce}
x^{\CE}_{\rho}=\frac{\ln\!\left(V^{\rho}/V^{0}\right)}{\gamma(1-\nu)},
\qquad
V^{0}\equiv\frac{1}{I^{s}}\sum_{i=1}^{I^{s}}V^{i,0}_{30},
\qquad
V^{\rho}\equiv\frac{1}{I^{s}}\sum_{i=1}^{I^{s}}V^{i,\rho}_{30}.
\end{equation}
With $\nu>1$ (\autoref{tab:est}) both averages are negative, so the ratio $V^{\rho}/V^{0}$ is positive; a welfare-improving reform makes $V^{\rho}$ less negative, the ratio lies below one, and dividing its negative logarithm by $\gamma(1-\nu)<0$ gives $x^{\CE}_{\rho}>0$.
% src: tables/estimated_parameters.tex:11 (nu = 1.7029 > 1)
The WTP column of every reform table reports $100\,x^{\CE}_{\rho}$, the CE change in percent, the cohort's willingness to pay (WTP) for the reform as a share of consumption.
% src: 7_policy_reforms.tex:73 (WTP column = CE welfare change in percent)

The leisure-equivalent (LE) change scales leisure instead of consumption:
\begin{equation}\label{eq:le}
x^{\LE}_{\rho}=\frac{\ln\!\left(V^{\rho}/V^{0}\right)}{(1-\gamma)(1-\nu)}=\frac{\gamma}{1-\gamma}\,x^{\CE}_{\rho}.
\end{equation}
% src: 12_app_quant.tex:582 (LE formula)
At the estimate $\gamma=0.5354$ the factor $\gamma/(1-\gamma)$ is 1.15, so the LE column of \autoref{tab:elim} is a rescaling of the CE column and appears nowhere else.
% src: tables/estimated_parameters.tex:9 (gamma = 0.5354); 1.15 = gamma/(1-gamma) is the identity stated in BLUEPRINT section 4.2, eq. 21 (not a result)

The dollar value of a reform converts $x^{\CE}_{\rho}$ with baseline consumption:
\begin{equation}\label{eq:wtp}
\mathcal{W}_{\rho}=x^{\CE}_{\rho}\,\Cbar{s},
\qquad
\Cbar{s}\equiv\sum_{t=s}^{85}(1+r)^{-(t-62)}\,\E\!\left[c^{0}_{t}\right],
\qquad s\in\{62,30\},
\end{equation}
% src: 12_app_quant.tex:632 (CE change times the NPV of per capita baseline consumption from 62, discounted to 62); 12_app_quant.tex:639 (from 30 for K); 5_quantitative_model.tex:23 (terminal age 85)
%TODO-AUTHORS: confirm that E[.] averages over all N simulated individuals with zero consumption after death (per-capita NPV); the sentence below changes if only survivors are averaged.
where $\E[\cdot]$ averages over the simulated cohort, $r$ is the net return of \autoref{sec:model}, and $\Cbar{s}$ is the per-capita net present value (NPV) at age 62 of status-quo consumption from age $s$ onward. The window is $s=62$ for the Social Security, Medicare, and income-tax families, whose reforms change flows received from age 62 onward, and $s=30$ for the saving subsidies, which change the return on assets at every age; for $t<62$ the factor compounds forward.
% src: 12_app_quant.tex:632, 639; tables/parameters_ss.tex:46 (ERA 62)
The CE change scales consumption at all ages while $\Cbar{62}$ counts consumption from age 62, so $\mathcal{W}_{\rho}$ is the dollar value of the CE change applied to consumption after 62; because $\Cbar{30}$ and $\Cbar{62}$ differ, MVPF levels of the saving-subsidy family are not compared with those of the other families.
% src: 12_app_quant.tex:632 (consumption window from 62), 639 (from 30 for K)

\subsection{Direct, indirect, and net fiscal effects}\label{sec:fiscal}

Three aggregates are computed for each family, all per capita and discounted to age 62 at the net return $r=3.2$ percent with the factor $(1+r)^{-(t-62)}$: $B$, the outlay of the program being reformed; $O$, the other taxes and transfers listed in \autoref{tab:conventions}; and $G$, government spending on the cohort, whose baseline value $G^{0}$ is the base of the Net column and whose change satisfies $\Delta G=\Delta B+\Delta O$ when every flow the reform changes is counted in $B$ or in $O$.
% src: 5_quantitative_model.tex:390 (net return 3.2 percent); 12_app_quant.tex:611 (discounted to age 62)
\autoref{tab:conventions} states what $B$ and $O$ contain in each family. The three columns are
\begin{equation}\label{eq:fiscal}
\text{Direct}_{\rho}=100\,\frac{\Delta B}{\text{base}_f},
\qquad
\text{Indirect}_{\rho}=100\,\frac{\Delta O}{O^{0}},
\qquad
\text{Net}_{\rho}=100\,\frac{\Delta G}{G^{0}},
\qquad \Delta G=\Delta B+\Delta O,
\end{equation}
where $\text{base}_f=B^{0}$ (baseline retirement benefits) for the Social Security families, $T^{p,0}$ (baseline payroll-tax collections) for Medicare, $|T^{0}|$ (baseline income-tax collections) for the income-tax reforms, for which $B$ is revenue forgone, and $G^{0}$ for the saving subsidies, for which $B$ is the subsidy. The Direct bases are positive by construction, so a negative Direct entry is a reduction in the program's own outlay, and a negative Net entry is a reduction in government spending on the cohort; the Indirect column is read with the same convention, a negative entry being a reduction in the government's other net outlays.
% src: 12_app_quant.tex:611, 615, 620-622, 626 (bases by family); 12_app_quant.tex:616 (Net base: NPV of government spending at the baseline)
% derived check, not printed (Open question 4): on the printed Social Security cells Net = 0.137 Direct + 1.12 Indirect to within 0.01 in every row
%   (e.g. AN3: 0.137(-16.75) + 1.12(-0.08) = -2.38; P1.2: 0.137(1.23) + 1.12(-0.07) = 0.09), so omega_D, omega_O > 0 and B^0 + O^0 differs from G^0;
%   the sign reading of the Indirect column rests on O^0 > 0, which these weights imply but the code's definitions have not confirmed.
%TODO-AUTHORS: Open question 6 -- confirm that the printed Direct column for I1-I5 is lost revenue over baseline income-tax collections (12_app_quant.tex:620 says "percentage change in the NPV of negative total income tax collection"; the printed cells are positive). Open question 4 -- confirm O^0 > 0 and G^0 > 0 under the code's definitions.

Because the three bases differ, the Net column is not the sum of the other two: for I1 the Direct and Indirect entries are 0.70 and $-0.81$ while Net is 0.69; for P1.1 they are 0.57 and $-0.10$ while Net is $-0.03$.
% src: tables/welfare_tax_mvpf_v8.tex:114 (I1: 0.70, -0.81, 0.69); tables/welfare_progressivity_mvpfv7.tex:115 (P1.1: 0.57, -0.10, -0.03)
The relation that holds is
\begin{equation}\label{eq:netidentity}
\text{Net}_{\rho}=\omega_{D}\,\text{Direct}_{\rho}+\omega_{O}\,\text{Indirect}_{\rho},
\qquad
\omega_{D}\equiv\frac{\text{base}_f}{G^{0}},
\qquad
\omega_{O}\equiv\frac{O^{0}}{G^{0}},
\end{equation}
which follows from \eqref{eq:fiscal} whenever $\Delta G=\Delta B+\Delta O$. The weights are ratios of baseline aggregates and are not printed; they sum to one only if $G^{0}=B^{0}+O^{0}$, which is not assumed. For the saving subsidies both equal one (the subsidy is zero at the baseline and the Direct base is $G^{0}$), and Net equals Direct plus Indirect in every K row.
% src: tables/welfare_savings_mvpf_v3.tex:120-123 (K1: 1.77 - 1.69 = 0.08; K2: 3.77 - 3.59 = 0.18; K3: 8.44 - 7.10 = 1.34)
%TODO-AUTHORS: Open question 4 -- the printed Social Security cells satisfy (eq:netidentity) with constant weights that do not sum to one (derived check above), so the code's G^0 is not B^0 + O^0 with O as described at 12_app_quant.tex:611; the items of G^0 must be stated from the code before the Net-base cell of Table 4 can say more than 'baseline government spending'.

% tables/new_accounting_conventions.tex -- Table 4 (tab:conventions), written by W2.
% DEFINITIONAL table (Override O1): no numeric constants. Every cell paraphrases an outcome definition of the
% May 2025 appendix; the notes quote the four family definitions verbatim.
% Sources of every definition in this file:
%   Social Security families (P, L, AN, AD, AL): frot/sub2025/12_app_quant.tex:611; 7_policy_reforms.tex:73 (NPV at 62 footnote)
%   Medicare Part A / Part B:                    frot/sub2025/12_app_quant.tex:615
%   Income tax I1-I5:                            frot/sub2025/12_app_quant.tex:620-622
%   Saving subsidies K1-K3:                      frot/sub2025/12_app_quant.tex:626
%   MVPF, dollar conversion, lower bound:        frot/sub2025/12_app_quant.tex:631-634 (s = 62), 12_app_quant.tex:639 (s = 30 for K)
%   Upper bound for K:                           frot/sub2025/7_policy_reforms.tex:423
%   MVPF cell forms:                             tables/welfare_progressivity_mvpfv7.tex:115 ([LB,+inf)); tables/welfare_age_dependence_mvpfv5.tex:108-126 (number for AN, inf for AD/AL); tables/welfare_savings_mvpf_v3.tex:120-123 ([LB,UB])
%   Discounting at the net return r:             frot/sub2025/5_quantitative_model.tex:390; 12_app_quant.tex:611
%TODO-AUTHORS: Open question 4 (BLUEPRINT section 11) -- the Net base is the "NPV of government spending at the baseline"
%  (12_app_quant.tex:616), whose items the source does not state; the printed cells rule out B^0 + O^0 with O as described at
%  12:611 (see the derived check in sections/05_accounting.tex). Confirm the items and sign convention of G^0 from the code;
%  no cell changes unless the code's O or G differs from 12_app_quant.tex:611-626.
\begin{table}[!t]
\centering
\caption{Accounting Conventions by Reform Family}
\label{tab:conventions}
\vspace{6pt}
\scriptsize
\setlength{\tabcolsep}{3pt}
\renewcommand{\arraystretch}{1.0}
\begin{threeparttable}
\begin{tabular}{@{}>{\raggedright\arraybackslash}p{0.85in}>{\raggedright\arraybackslash}p{1.05in}>{\raggedright\arraybackslash}p{0.8in}>{\raggedright\arraybackslash}p{0.9in}>{\raggedright\arraybackslash}p{0.7in}>{\raggedright\arraybackslash}p{0.55in}>{\raggedright\arraybackslash}p{1.15in}@{}}
\toprule
Family & Program outlay $B$ & Base of Direct ($\text{base}_f$) & Base of Indirect ($O^{0}$) & Base of Net ($G^{0}$) & Window for $\mathcal{W}$ ($\Cbar{s}$) & MVPF cell \\
\midrule
Social Security: benefit formula (P, L); claiming schedule (AN, AD, AL)
 & NPV of retirement benefits paid by the program, $b_t\,ss_t$, from age 62
 & $B^{0}$: baseline retirement benefits
 & All other taxes and transfers: disability benefits and transfers, net of income and payroll taxes
 & Baseline government spending
 & From age 62
 & Number if $\Delta G>0$; $[\,\mathcal{W}/\Delta B,+\infty)$ if $\Delta G<0<\Delta B$; $\infty$ if $\Delta G<0$ and $\Delta B<0$; for $\mathcal{W}<0$, $\Delta G<0$ (AN) the number is the value of the reverse reform \\
\addlinespace
Medicare Part A; Part B
 & NPV of medical payments by the program to the targeted beneficiaries after age 65
 & $T^{p,0}$: baseline payroll-tax collections
 & All taxes and transfers: Social Security benefits and transfers, net of income and payroll taxes
 & Baseline government spending
 & From age 62
 & Number \\
\addlinespace
Income tax (I1--I5)
 & NPV of income-tax revenue forgone after age 62, $-\Delta T$
 & $|T^{0}|$: baseline income-tax collections
 & All other taxes and transfers: Social Security benefits and transfers, net of payroll taxes
 & Baseline government spending
 & From age 62
 & Number \\
\addlinespace
Saving subsidy (K1--K3)
 & NPV of the subsidy $\mathit{Sub}$ paid on assets at all ages
 & $G^{0}$: baseline government spending
 & Government spending excluding the subsidy
 & Baseline spending including the subsidy
 & From age 30, compounded to 62
 & $[\,\mathcal{W}/\mathit{Sub},$ $\mathcal{W}/\Delta G^{\mathrm{full}}]$ \\
\bottomrule
\end{tabular}
\begin{tablenotes}[flushleft]\scriptsize
\item[] \textit{Notes:} All fiscal aggregates are per-capita net present values (NPVs) discounted to age 62 at the net return $r$ with the factor $(1+r)^{-(t-62)}$; the columns are defined in \eqref{eq:ce}, \eqref{eq:wtp}--\eqref{eq:netidentity}, and \eqref{eq:mvpf}--\eqref{eq:mvpfk}, and the three fiscal columns are related by \eqref{eq:netidentity}, not by addition. Participation columns report changes in labor force participation at ages 30--49, 62--64, and 65--69 in percentage points. The definitions used in computing the tables are the following.
% src: 7_policy_reforms.tex:73 (participation age groups)
\emph{Social Security:} ``the direct fiscal effect is calculated as the percentage change in the net present value (NPV) of total retirement benefits paid out by the program, discounted to age 62 by a factor of $(1+r)^{-1}$. The indirect fiscal effects are computed as the percentage change in the NPV of all other taxes and transfers, including payments to disability benefits and government transfers net of income and payroll tax collections. The net fiscal effect combines the direct and indirect effects.'' The factor $(1+r)^{-1}$ in this definition is applied for each year between $t$ and age 62, that is, $(1+r)^{-(t-62)}$.
% src: 12_app_quant.tex:611; 7_policy_reforms.tex:192 ("(1+r)^{-1(t-62)}")
\emph{Medicare:} ``the direct fiscal effect is calculated as the changes in the NPV of medical payments made by the program to the targeted beneficiaries after age 65 (Part A for everyone; Part B for labor market nonparticipants), expressed as a share of payroll tax collections at the baseline. The indirect fiscal effects are computed as the percentage change in the NPV of all taxes and transfers, including payments to all Social Security benefits and government transfers net of income and payroll tax collections. The net fiscal effect is measured as the sum of the NPV of the change in government spending and the total medical payments to the targeted beneficiaries, divided by the NPV of government spending at the baseline.''
% src: 12_app_quant.tex:615
\emph{Income tax:} ``the direct fiscal effect is calculated as the percentage change in the NPV of negative total income tax collection. The indirect fiscal effect is computed as the percentage change in the NPV of all other taxes and transfers, including payments to all Social Security benefits and government transfers net of payroll tax collections. The net fiscal effect combines the direct and indirect effects.''
% src: 12_app_quant.tex:620-622
\emph{Saving subsidy:} ``the direct fiscal effect is calculated as the NPV of the subsidy as a share of government spending at the baseline. The indirect fiscal effect is computed as the percentage change in the NPV of government spending, excluding subsidies. The net fiscal effect is measure[d] as the percentage change in the NPV of the government spending, including subsidies.'' The welfare NPV for this family starts at age 30 and ``the lower bound uses only the subsidy in the denominator.''
% src: 12_app_quant.tex:626; 12_app_quant.tex:639
\end{tablenotes}
\end{threeparttable}
\end{table}


\subsection{The marginal value of public funds and the table columns}\label{sec:mvpf}

The marginal value of public funds (MVPF) of a reform is the ratio of the beneficiaries' willingness to pay to the net cost to the government, inclusive of behavioral responses \citep{hendren2016policy,hendren2020unified,finkelstein2020welfare}. Here both terms are per capita for the cohort:
\begin{equation}\label{eq:mvpf}
\MVPF_{\rho}=\frac{\mathcal{W}_{\rho}}{\Delta G_{\rho}}.
\end{equation}
% src: 12_app_quant.tex:631-632
Substituting \eqref{eq:wtp} and \eqref{eq:fiscal},
\begin{equation}\label{eq:mvpfcolumns}
\MVPF_{\rho}=\kappa_f\,\frac{\WTP_{\rho}}{\text{Net}_{\rho}},
\qquad
\kappa_f\equiv\frac{\Cbar{s}}{G^{0}},
\end{equation}
so within a family the MVPF column is proportional to WTP/Net. For P1.2, the ratio of WTP to Net is 0.22/0.09; the MVPF of 0.97 converts it to dollars of welfare per dollar of net spending.
% src: tables/welfare_progressivity_mvpfv7.tex:116 (P1.2: MVPF 0.97, WTP 0.22, Net 0.09)

When $\mathcal{W}_{\rho}>0$ and $\Delta G_{\rho}>0$ the MVPF is a positive number. When $\mathcal{W}_{\rho}>0>\Delta G_{\rho}$ the reform raises welfare and reduces net spending: it is self-financing, the MVPF is infinite, and the tables print $\infty$; when in addition $\Delta B_{\rho}>0$ they report a lower bound that counts only the program's own outlay,
\begin{equation}\label{eq:mvpflb}
\MVPF^{\mathrm{LB}}_{\rho}=\frac{\mathcal{W}_{\rho}}{\Delta B_{\rho}}=\frac{\Cbar{s}}{\text{base}_f}\,\frac{\WTP_{\rho}}{\text{Direct}_{\rho}},
\end{equation}
% src: 12_app_quant.tex:634 (lower bound = dollar CE gain over the change in the direct cost)
printed as $[\MVPF^{\mathrm{LB}}_{\rho},+\infty)$; P1.1 is the one such row, with a lower bound of 0.52.
% src: tables/welfare_progressivity_mvpfv7.tex:115 ([0.52,+infty))
When $\Delta B_{\rho}<0$ as well as $\Delta G_{\rho}<0$, as in every AD and AL row, no finite bound exists.
% src: tables/welfare_age_dependence_mvpfv5.tex:112-126 (WTP > 0, Direct < 0, Net < 0, MVPF = infty in all AD and AL rows)
When $\mathcal{W}_{\rho}<0$ and $\Delta G_{\rho}<0$, as in the AN rows, the positive ratio is the welfare loss per dollar of net savings and, since $x^{\CE}$ and $\Delta G$ both change sign when a reform is reversed, equals the MVPF of the reverse reform except for the difference between the consumption aggregates $\Cbar{s}$ of \eqref{eq:wtp} under the two rules; \autoref{sec:claiming} reads the AN rows in this way.
% src: tables/welfare_age_dependence_mvpfv5.tex:108-110 (AN1-AN3: WTP < 0, Net < 0, finite MVPF); 7_policy_reforms.tex:119 (reversal reading)
For the saving subsidies the tables print an interval,
\begin{equation}\label{eq:mvpfk}
\MVPF_{\rho}\in\left[\frac{\mathcal{W}_{\rho}}{\mathit{Sub}_{\rho}},\;\frac{\mathcal{W}_{\rho}}{\Delta G^{\mathrm{full}}_{\rho}}\right],
\end{equation}
whose upper bound uses the net cost $\Delta G^{\mathrm{full}}_{\rho}$ computed as if all capital gains were realized and taxed in each period.
% src: 12_app_quant.tex:639 (lower bound uses only the subsidy); 7_policy_reforms.tex:423 (upper bound: all capital gains realized and taxed)

\citet{hendren2020unified} write that ``the creation of MVPFs requires various judgment calls.'' Ours are the following.
% src: 7_policy_reforms.tex:73 (quotation from hendren2020unified)
Welfare is the equal-weighted ex ante average of expected utility at age 30. Fiscal flows are per-capita NPVs at age 62 at the net return $r$ and enter at face value: no marginal cost of public funds multiplies the Net column, unlike in the marginal-cost-of-funds tradition of \citet{mayshar1990measures}, \citet{slemrod1996costs}, and \citet{kleven2006marginal}; the reform's own revenue response is inside $\Delta G$, and the cost of raising the funds enters by comparing MVPFs across policies. The accounting is for one cohort, with no transition (\autoref{sec:institutions} states the stance on incidence outside the cohort), and at fixed prices (\autoref{sec:scope} sets out how factor-price responses would enter $\mathcal{W}$ and $\Delta G$). The participation columns report changes in labor force participation at ages 30--49, 62--64, and 65--69 in percentage points (\pp); they are the behavioral responses behind the fiscal effects and enter no formula directly.
% src: 7_policy_reforms.tex:73 (participation age groups, in percentage points)

\subsection{Relation to sufficient statistics: consumption smoothing and the fiscal multiplier}\label{sec:suffstats}

In the sufficient-statistics approach to social insurance \citep{baily1978some,chetty2006general,chetty2008moral,chetty2009sufficient,chetty2013social}, a marginal expansion of a program raises welfare when the consumption-smoothing value of the transfer exceeds its fiscal cost, the mechanical outlay plus the fiscal externality from behavioral responses. \citet{haller2022welfare} extends the formula to a dynamic retirement setting and evaluates increases in the early retirement age and reductions in benefit levels with a fiscal multiplier, the total fiscal effect per mechanical dollar, estimated from Austrian administrative data; the multiplier is the object to which the fiscal columns are mapped below.
%TODO-AUTHORS: Open question 16 -- haller2022welfare is an @unpublished entry; the instruments named above (early retirement age; pension generosity) are those of the 2019 CESifo draft and the 2022 WIFO seminar abstract; confirm against the current version before this sentence is final. No number is imported from it.
Split the change in program outlays into a mechanical and a behavioral part,
\begin{equation}\label{eq:decomp}
\Delta B=\Delta B^{M}+\Delta B^{B},
\end{equation}
where $\Delta B^{M}$ applies the reformed benefit, coverage, tax, or return rule to the status-quo histories with every policy function held at its status-quo value, and $\Delta B^{B}$ is the remainder; $\Delta O$ is behavioral except for mechanical interactions through the consumption floor \eqref{eq:transfer} and the income-tax base in \eqref{eq:budget}. Define
\begin{equation}\label{eq:fm}
\FM\equiv\frac{\Delta G}{\Delta B^{M}}=1+\frac{\Delta B^{B}+\Delta O}{\Delta B^{M}},
\qquad
\FM_{B}\equiv\frac{\Delta B}{\Delta B^{M}}=1+\frac{\Delta B^{B}}{\Delta B^{M}},
\end{equation}
the fiscal multiplier and its within-program part, and
\begin{equation}\label{eq:bcmap}
\hat{\pi}\equiv\frac{\mathcal{W}}{\Delta B^{M}}-1,
\qquad
\MVPF=\frac{1+\hat{\pi}}{\FM},
\qquad
\MVPF^{\mathrm{LB}}=\frac{1+\hat{\pi}}{\FM_{B}}.
\end{equation}
Both equalities in \eqref{eq:bcmap} are identities: $\mathcal{W}=(1+\hat\pi)\,\Delta B^{M}$ by the definition of $\hat\pi$ and $\Delta G=\FM\,\Delta B^{M}$ by \eqref{eq:fm}, so their ratio is \eqref{eq:mvpf}, and replacing $\Delta G$ by $\Delta B$ gives \eqref{eq:mvpflb}; both require only $\Delta B^{M}\neq0$ and hold exactly for discrete reforms. An MVPF of one is the condition $1+\hat\pi=\FM$, the Baily--Chetty condition when a dollar of public funds is valued at one, that is, relative to a financing instrument whose own MVPF is one; an infinite MVPF with $\Delta B^{M}>0$ means $\FM<0$, behavioral fiscal gains exceeding the mechanical outlay.

For a marginal change in a transfer rule, the envelope theorem applied to \eqref{eq:bellman} gives $\mathcal{W}\approx\sum_{t}(1+r)^{-(t-62)}\,\E[\omega_t\,\delta^{M}_t]$, where $\delta^{M}_t$ is the mechanical change in the individual's resources at age $t$ and $\omega_t$ is the marginal utility of consumption at $t$, discounted at $\beta$ and compounded to age 62 at $r$, relative to the average marginal utility implicit in the conversion of $x^{\CE}$ to dollars (\autoref{app:derivations}). Then $\pi\approx\E[\omega\,\delta^{M}]/\E[\delta^{M}]-1$ is the transfer-weighted average of relative marginal utilities minus one, the consumption-smoothing value of the transfer: the term of \citet{baily1978some} and \citet{chetty2006general}, positive when the transfer accrues in states of high marginal utility such as poor health, low assets, and survival to old age \citep[see also][]{landais2021value,kolsrud2024retirement}. With nonseparable utility and $\nu>1$, marginal utility of consumption falls with leisure, so a dollar is valued less after labor-force exit than while working at the same consumption.
% src: tables/estimated_parameters.tex:11 (nu > 1); eq:utility (u_c proportional to l^{(1-gamma)(1-nu)}, decreasing in l when nu > 1)
Two cautions apply. First, the structural $\mathcal{W}$ is exact for the discrete reforms considered, whereas the envelope expression is first-order in the size of the reform and omits the utility changes of individuals who switch participation or claiming state, discrete choices with fixed costs \citep{kleven2021sufficient}; $\hat\pi$ is defined by \eqref{eq:bcmap} so that the identities hold regardless. Second, the printed Direct column is $\Delta B$, not $\Delta B^{M}$: it combines the mechanical change with the within-program behavioral response, while the Indirect column holds the behavioral response outside the program. Separating $\Delta B^{M}$ from $\Delta B^{B}$ requires an accounting pass over saved status-quo histories with behavior held fixed, which is not part of the present output: the tables identify $(1+\hat\pi)/\FM$ and $(1+\hat\pi)/\FM_B$ but neither $\FM$ nor $\hat\pi$ separately, the fiscal multiplier cannot be read from the Direct column, and the comparison with \citet{haller2022welfare} is one of objects, not of magnitudes.
%TODO-AUTHORS: Open question 4 -- if Delta B^M can be computed by an accounting pass over saved panels, FM and pi-hat become numbers and this paragraph is upgraded.

Two sign inferences follow from the printed cells and the schedules of \autoref{tab:schedules}, and no others. For AD1--AD5 every factor after the NRA rises and every factor before it is unchanged, so with claiming held at its status-quo age outlays rise for everyone who claims after age 65 or whose benefits withheld under the earnings test after the NRA are credited at the DRC (\autoref{subsec:programs}), and fall for no one: $\Delta B^{M}\ge0$.
% src: tables/DRC_visualization_4.tex:37 (AN0) versus :43-47 (AD1-AD5: factors at 62-65 unchanged, factors at 66-70 higher); tables/parameters_ss.tex:47 (NRA 65)
Direct and Net are negative in all five rows, so $\Delta B<0$ and $\Delta G<0$, hence $\FM_B<0$ and $\FM<0$ whenever $\Delta B^{M}>0$: the behavioral change in outlays more than offsets the mechanical increase.
% src: tables/welfare_age_dependence_mvpfv5.tex:112-116 (AD1-AD5: Direct -2.46 to -5.05, Net -0.45 to -0.78)
For AN1--AN3 every factor falls, so $\Delta B^{M}<0$ by construction, and WTP is negative; the finite ratio is the reversal value of \autoref{sec:mvpf}.
% src: tables/DRC_visualization_4.tex:38-40 (AN1-AN3: every factor below the AN0 factor at the same age); tables/welfare_age_dependence_mvpfv5.tex:108-110 (WTP -0.38, -0.74, -1.09)
For the benefit-formula, Medicare, and income-tax families $\Delta B^{M}>0$, but the sign of $\Delta B^{B}$ is not identified, so whether $\FM$ exceeds one is undetermined; the AL rows move factors in both directions, so $\Delta B^{M}$ itself is unsigned.
% src: tables/DRC_visualization_4.tex:50-58 (AL1-AL9: factors below the NRA fall, factors above it rise)
